\chapter{Допуск полного шумового кокасательного носителя}
\label{ch:v8-full-noise-cotangent-carrier-admission-gate}

Полное калибровочное замыкание содержит $15$ комплексных направлений переноса и
$12$ эрмитовых калибровочных направлений. Их нельзя складывать как однородное
комплексное пространство: калибровочные коэффициенты вещественны. Правильная
овеществление даёт
\begin{equation}
 \dim_{\mathbb R}\mathcal N_{\rm mix}=2\cdot15+12=42,
\end{equation}
что точно совпадает с $42$-мерным вещественным пространством полей полной
суперсвязности. На уровне типов полный кокасательный родитель поэтому допустим.

Если ошибочно считать все $15+12=27$ направлений комплексными, получается
$54$ вещественных координаты и ложный размерностный конфликт.

Однако текущий примитивный QMS использует только $25$ самосопряжённых
операторов скачка: один связывающий, двенадцать калибровочных и двенадцать межсекторных. Поэтому
его кадр имеет коразмерность
\begin{equation}
 42-25=17
\end{equation}
в полном смешанно-вещественном носителе. Он не может быть невырожденным
кокасательным репером всего пространства полей без добавления отсутствующих направлений переноса
направлений и повторного вычисления общей следовой метрики.

\begin{theorem}[Типизированный допуск полного носителя]
Смешанный вещественный шумовой модуль
$\operatorname{Realify}(\mathbb C^{15})\oplus\mathbb R^{12}$ имеет
размерность $42$ и допускает спаривание с полным пространством полей.
\end{theorem}

\begin{nogocriterion}[Недостаточность текущего 25-скачковый процесса]
Нельзя объявлять существующий полный QMS кокасательной динамикой всего
42-мерного полевого пространства: его реперу скачков недостаёт 17 вещественных
направлений. Нельзя также устранять разрыв формальной комплексфикацией
эрмитовой калибровочной алгебры, поскольку она удваивает физические координаты.
\end{nogocriterion}

\begin{protocol}
\begin{enumerate}
 \item часть переноса: \textbf{$15$ комплексных измерений, $30$ вещественных измерений};
 \item калибровочная часть: \textbf{$12$ самосопряжённых вещественных направлений};
 \item полный смешанный носитель: \textbf{$42$ вещественных измерений};
 \item пространство полей суперсвязностей: \textbf{$42$ вещественных измерения};
 \item наивная комплексфикация: \textbf{$54$ вещественных измерений, отвергнута};
 \item текущий репер скачков: \textbf{$25$ вещественных измерений};
 \item недостающих направлений: \textbf{$17$};
 \item следующий гейт: \textbf{построение полного 42-мерного репера скачков и
       его общей метрик следаи}.
\end{enumerate}
\end{protocol}

Машинный сертификат сохранён в
\path{s2t_v8_full_noise_cotangent_carrier_admission_gate_results.json}.